% Sample-size sensitivity of NMI under a null where all classes are i.i.d.
% Source of record: theoretical null-model simulation (T=1000 trials per config).
% Supports the clustering cross-check (Section~\ref{sec:eval-clustering}): explains why
% small-sample NMI can be elevated even with no class-distinguishing signal.
\subsection{Sample-size sensitivity of NMI}
\label{sec:eval-nmi-sensitivity}

The clustering cross-check (Section~\ref{sec:eval-clustering}) treats ARI as the primary
conservative indicator because NMI can take moderately high values under small samples.
To justify that choice, we quantify how \emph{finite-sample effects} alone can inflate NMI
even when there is \emph{no} discriminative signal, using an idealized null in which all
classes are statistically indistinguishable.

\paragraph{Null hypothesis.} We sample every class from the \emph{same} distribution,
$x_{i}\sim\mathcal{N}(0,I_{d})$ for all classes. Since no class carries class-specific
information, a clustering algorithm should not recover class membership, and both NMI and
ARI should concentrate around $0$ as the sample size grows.

\paragraph{Simulation.} For each configuration we run $T=1000$ independent trials. In each
trial we (i)~generate $C$ classes with $n$ i.i.d.\ samples per class from
$\mathcal{N}(0,1)$, (ii)~apply K-Means with $K=C$ to the pooled data, and (iii)~score the
predicted assignments against the ground-truth labels with NMI and ARI. We sweep $n$ while
holding $C$ fixed, and report the trial means $\overline{\mathrm{NMI}}$ and
$\overline{\mathrm{ARI}}$. We consider a \textbf{PI-like} binary structure ($C=2$,
$n\in\{3,10,50,100\}$) and an \textbf{ICD10-like} multi-class structure ($C=5$,
$n\in\{3,10,100\}$).

\paragraph{Results.} Table~\ref{tab:nmi-sample-size-sensitivity} shows that, although all
classes are indistinguishable by construction, \textbf{NMI is noticeably positive at very
small $n$ but collapses toward zero as $n$ grows}, whereas \textbf{ARI stays near zero
throughout}. For the binary structure, $\overline{\mathrm{NMI}}$ falls from $0.1808$ at
$n{=}3$ ($N{=}6$) to $0.0404$ at $n{=}10$, and to $0.0066$/$0.0038$ at $n{=}50$/$100$,
while $\overline{\mathrm{ARI}}$ stays at $-0.0190$, $-0.0020$, $-0.0011$, $0.0002$. For the
multi-class structure the same trend holds: $\overline{\mathrm{NMI}}$ goes from $0.4168$
($n{=}3$) to $0.1249$ ($n{=}10$) to $0.0105$ ($n{=}100$), with $\overline{\mathrm{ARI}}$
near zero ($0.0005$, $0.0018$, $-0.0000$). Hence moderately high NMI in small-sample
clustering should not be read as actionable leakage, and ARI is the more conservative
indicator of practical recoverability in our unsupervised evaluation.

\begin{table}[ht]
  \centering
  \caption{Sample-size sensitivity of NMI and ARI under the null hypothesis where all
    classes are i.i.d.\ sampled from $\mathcal{N}(0,1)$ ($T{=}1000$ trials per row).}
  \label{tab:nmi-sample-size-sensitivity}
  \begin{tabular}{lrrrrr}
    \toprule
    Structure & $n$ (per class) & $C$ & $N$ & $\overline{\mathrm{NMI}}$ & $\overline{\mathrm{ARI}}$ \\
    \midrule
    PI-like    & 3   & 2 & 6   & $0.1808$ & $-0.0190$ \\
    PI-like    & 10  & 2 & 20  & $0.0404$ & $-0.0020$ \\
    PI-like    & 50  & 2 & 100 & $0.0066$ & $-0.0011$ \\
    PI-like    & 100 & 2 & 200 & $0.0038$ & $0.0002$  \\
    \midrule
    ICD10-like & 3   & 5 & 15  & $0.4168$ & $0.0005$  \\
    ICD10-like & 10  & 5 & 50  & $0.1249$ & $0.0018$  \\
    ICD10-like & 100 & 5 & 500 & $0.0105$ & $-0.0000$ \\
    \bottomrule
  \end{tabular}
\end{table}
